% Данный файл распространяется под лицензией CC BY 4.0. Текст лицензии размещён на https://creativecommons.org/licenses/by/4.0/
% (c) Кафедра системного программирования, 2026

% Компилировать XeTeX

\documentclass[a4paper]{article}

\usepackage[a4paper, top=8mm, bottom=8mm, left=8mm, right=8mm]{geometry}

\usepackage{polyglossia}
\setdefaultlanguage[babelshorthands=true]{russian}
\setotherlanguage{english}

\usepackage{fontspec}
\setmainfont{FreeSerif}
\newfontfamily{\russianfonttt}[Scale=0.7]{DejaVuSansMono}

\usepackage[tiny, compact]{titlesec}

\usepackage{titling}
\setlength{\droptitle}{-1cm}
\pretitle{\begin{center}\begin{bfseries}\Large}
\posttitle{\par\end{bfseries}\end{center}}
\preauthor{\begin{center}\normalsize}
\postauthor{\par\end{center}\vspace{-1.8cm}}

\usepackage{hyperref}
\usepackage{bookmark}
\usepackage{csquotes}
% Сюда можно дописывать нужные вам пакеты.

\title{Название работы}

\author{ФИО}

% Дата много места занимает, её лучше не указывать.
\date{}

\begin{document}

\maketitle

\begin{flushright}
    Группа: \emph{группа в формате \enquote{<год поступления>.<Б/М><номер группы>-мм}}

    Кафедра: \emph{кафедра, на которой планируется защита/сдача работы, например, "системного программирования"}

    % Степень/звание и должность научника мы лучше вас знаем, их можно не указывать.
    Научный руководитель: \emph{ФИО научного руководителя}

    % А вот про консультанта укажите официальное место работы, с "ООО" или чем-то таким, и должностью желательно по трудовой, не просто "techlead". Выясните у консультанта.
    % Если консультанта нет (а на первом курсе у вас скорее всего его нет), просто удалите это.
    Консультант: \emph{место работы, должность, степень, звание (если есть) консультанта}
    
    Номер семестра практики: \emph{семестр, за который практика}
\end{flushright} 

\section{Введение}

Пара абзацев про предметную область и зачем вы делаете то, что делаете, кому это может быть полезно. Только без воды, никакого \enquote{в современном мире}! И без фантазий на тему продвижения границ человеческого знания, пишите как есть~---например, это учебная задача, которая позволит освоить то-то и то-то и оставить будущим поколениям такой-то пример. Если практика выполняется в рамках какого-то проекта или в компании, укажите это явно, с названием проекта или компании и кратким описанием, что это за проект.

\textbf{Требуемый размер всего текста практики~--- три страницы}, так что не увлекайтесь. Учитесь писать максимально сжато, чётко, по сути и с фактами~--- в эпоху больших языковых моделей полотна текста и красивый литературный стиль ценности не имеют.

\section{Постановка задачи}

\textbf{Целью работы является} \dots~--- тут кратко сформулируйте цель проекта, т.е. что конкретно вы делаете (зачем~--- не надо, это должно быть во введении).

\textbf{Для достижения цели требуется решить следующие задачи:}
\begin{enumerate}
    \item Выполнить обзор существующих решений.
    \item Спроектировать решение.
    \item Реализовать решение.
    \item Выполнить апробацию и/или эксперименты.
\end{enumerate}

Задачи --- это по сути план работы. Обязательно должна быть задача про обзор чего-то (предметной области, алгоритма, существующих решений, проекта, в который вы встраиваетесь и т.п.), чтобы было видно, что вы не изобретали велосипед. Задачи на проектирование и реализацию для небольших семестровых практик первого курса можно не разделять, но если там пришлось попроектировать, вынесите это отдельным результатом. Обязательна задача на проверку того, что получилось~--- это может быть апробация (т.е. внешняя валидация, дали другим людям и собрали обратную связь, или выступили на конференции и собрали обратную связь, или что-то такое) или эксперименты (запустили и что-то померили) или, в крайнем случае, тестирование (внутренняя верификация, вы проверили, что ваша реализация удовлетворяет вашим представлениям о желаемом результате). Лучше, когда есть одновременно всё перечисленное, но кратко.

\section{Обзор}

Обзор может быть обзором существующих решений, реализуемого алгоритма, используемых технологий, публикаций, решения, в которое вы встраиваетесь, предыдущих достижений исследовательской группы и т.п. Обзор должен иметь какую-то цель, которую желательно явно сформулировать в начале обзора (например, показать актуальность темы или выбрать хорошие решения для переиспользования). Всё, что сделано по работе не лично вами~--- должно быть в обзоре.

Однако обзор должен быть кратким~--- БЯМ накопают читателю сколько угодно деталей, если ему будет надо.

Обязательно сделайте краткие выводы из обзора, буквально одно-два предложения, например \enquote{Таким образом, решения, удовлетворяющего всем требованиям, не найдено, поэтому задача актуальна}.

\emph{Задача со звёздочкой (не требуется на первом курсе, но часть инженерной культуры)~--- явно указать критерии поиска и отбора обозреваемого материала, критерии сравнения, если что-то с чем-то сравнивается, и пояснить, почему они важны, представить в конце сводную таблицу с результатами сравнения по критериям}.

\section{Описание предлагаемого решения}

Тут желательно осветить общую структуру решения (т.е. из каких компонентов оно состоит и какой компонент за что отвечает), используемые технологии (с указанием конкретных версий), интересные технические детали (\enquote{Всё, над чем пришлось думать больше пяти минут}, как говорил Андрей Николаевич Терехов). По возможности обоснуйте все принятые решения. Не увлекайтесь деталями, у вас всё равно будет ссылка на код в заключении, там в комментариях или README можно пояснить всё, что очень хотелось рассказать тут, но не влезло в три страницы.

\section{Апробация / Эксперименты / Тестирование}

% Название секции можно поменять на более подходящее вашей работе

Название этой секции --- либо \enquote{Апробация}, либо \enquote{эксперименты}, либо \enquote{тестирование} (если содержательно есть и то и то, можно две отдельные секции сделать). Тут опишите, как проводилась проверка того, что получилось что-то адекватное. Если есть отзывы пользователей, отлично (особенно если они в структурированном виде, типа анкеты SUS). Если есть отзывы команды, консультанта и т.п., хорошо. Если есть только тесты. то хотя бы их тут опишите.

Если работа подразумевала замеры чего-то (например, производительности), приведите тут итоговые результаты и выводы, и дайте ссылку на таблицы с данными. Не забывайте про матстат (приводите матожидание и среднеквадратичное отклонение, не просто результат замера~--- вас пока не учили, поэтому придираться не будут, но опять-таки в эпоху БЯМ сделать хорошо можно и до прослушивания курса матстата).

Приведите выводы из эксперимента.

\section{Заключение}

В ходе выполнения работы получены следующие результаты:

% ниже списком перечисляете то, что выносится как результат практики.

\begin{itemize}
    \item \dots
    \item \dots
    \item \dots
\end{itemize}

Материалы, иллюстрирующие выполнение работы, размещены по адресу: \url{http://www.example.com} (если есть и применимо --- скриншоты, видео на 1-2 минуты и т.п. с демонстрацией; может быть очень полезно даже для консольных утилит или библиотек).

Репозиторий/ссылка на пуллреквест: \url{http://www.example.com} (это обязательно; ссылок может быть несколько).

Техническая документация: \url{http://www.example.com} (если применимо --- для пуллреквестов можно техническую сторону прямо в комментарии к пуллреквесту указать, но можно и в отдельную вики-страницу вынести).

\end{document}